\chapter{Annexes}
\label{chap:annexes}

\section{Schémas existants intégrés}
Le tableau~\ref{tab:schemas-integres} liste les schémas existants retenus dans le rapport.

\begin{table}[H]
\centering
\small
\caption{Schémas existants intégrés}
\label{tab:schemas-integres}
\begin{tabularx}{\textwidth}{p{3.3cm}p{3cm}Y p{3cm}}
\toprule
\textbf{Fichier} & \textbf{Emplacement} & \textbf{Légende} & \textbf{Objectif} \\
\midrule
\pathcode{figure_08.jpeg} & Chapitre~\ref{chap:conception} & Diagramme de cas d'utilisation du back-office gérant & Montrer le périmètre administratif. \\
\pathcode{figure_05.png} & Chapitre~\ref{chap:conception} & Diagramme de cas d'utilisation des modules salle et cuisine & Expliquer les interactions serveur/cuisine. \\
\pathcode{figure_02.png} & Chapitre~\ref{chap:conception} & Diagramme de cas d'utilisation du portail client & Montrer les fonctions client. \\
\pathcode{figure_03.png} & Chapitre~\ref{chap:conception} & Diagramme de classes de Tastify & Présenter les entités métier majeures en page paysage. \\
\pathcode{figure_06.png} & Page de garde & Logo Tastify & Identifier visuellement le projet. \\
\bottomrule
\end{tabularx}
\end{table}

\section{Schémas techniques créés dans le rapport}
Les schémas suivants ont été créés directement en LaTeX sous forme de diagrammes textuels. Ce ne sont pas des captures d'écran.

\begin{table}[H]
\centering
\small
\caption{Schémas techniques créés}
\label{tab:schemas-crees}
\begin{tabularx}{\textwidth}{p{3.5cm}p{3.8cm}Y}
\toprule
\textbf{Schéma} & \textbf{Emplacement} & \textbf{Contenu} \\
\midrule
Architecture fonctionnelle & Chapitre~\ref{chap:conception}, figure~\ref{fig:architecture-fonctionnelle} & Deux frontends, backend DRF, MySQL, Redis, Celery et Hugging Face. \\
Pipeline de sentiment & Chapitre~\ref{chap:sentiment}, figure~\ref{fig:pipeline-sentiment} & Création avis, tâche Celery, API Hugging Face, fallback local, stockage et dashboard. \\
Architecture technique & Chapitre~\ref{chap:realisation}, figure~\ref{fig:architecture-technique} & Ports, backend ASGI, base de données, Redis, worker et beat. \\
Séquence JWT & Chapitre~\ref{chap:realisation}, figure~\ref{fig:auth-jwt} & Connexion, vérification, access token, cookie refresh et refresh ultérieur. \\
Flux KDS & Chapitre~\ref{chap:realisation}, figure~\ref{fig:flux-kds} & Serveur, API commandes, base, Redis, KDS et cuisinier. \\
Déploiement Docker Compose & Chapitre~\ref{chap:realisation}, figure~\ref{fig:docker-compose} & Services Compose et dépendances principales. \\
\bottomrule
\end{tabularx}
\end{table}

\section{Captures intégrées dans le corps du rapport}
Les captures suivantes sont intégrées dans les chapitres de réalisation, d'analyse de sentiment et de tests. Elles ne sont plus présentées comme des éléments à ajouter ultérieurement.

\begin{table}[H]
\centering
\scriptsize
\caption{Captures placées dans le corps du rapport}
\label{tab:captures-integrees}
\begin{tabularx}{\textwidth}{p{3.8cm}p{3.2cm}Y}
\toprule
\textbf{Fichier} & \textbf{Section} & \textbf{Objectif} \\
\midrule
\pathcode{capture_login_staff.png} & Réalisation & Prouver la séparation staff et l'accès sécurisé. \\
\pathcode{capture_dashboard.png} & Réalisation & Montrer le pilotage et les indicateurs gérant. \\
\pathcode{capture_salle.png} & Réalisation & Montrer la gestion opérationnelle des tables. \\
\pathcode{capture_kds.png} & Réalisation & Montrer le suivi de préparation côté cuisine. \\
\pathcode{capture_client.png} & Réalisation & Montrer la consultation du menu côté client. \\
\pathcode{capture_stock_admin.png} & Réalisation & Illustrer le suivi des ingrédients et alertes. \\
\pathcode{capture_reservation_client.png} & Réalisation & Montrer la réservation côté client. \\
\pathcode{capture_checkout_payment.png} & Réalisation & Illustrer le paiement et le partage de l'addition. \\
\pathcode{capture_loyalty_client.png} & Réalisation & Montrer les points et récompenses. \\
\pathcode{capture_hr_admin.png} & Réalisation & Montrer les employés et la gestion RH. \\
\pathcode{capture_avis_sentiment.png} & Analyse de sentiment & Relier les avis clients au score de sentiment. \\
\pathcode{capture_test_pytest.png} & Tests & Prouver les résultats backend pytest. \\
\pathcode{capture_test_vitest.png} & Tests & Prouver les résultats Vitest des frontends. \\
\pathcode{capture_test_playwright.png} & Tests & Prouver les résultats Playwright client. \\
\pathcode{capture_test_ci.png} & Tests & Relier workflow CI et validation locale. \\
\bottomrule
\end{tabularx}
\end{table}

\section{Commandes utiles}
Les commandes ci-dessous servent de rappel pour vérifier le projet et le rapport.

{\scriptsize
\begin{verbatim}
cd docs/rapport-pfa-tastify-final
pdflatex --disable-installer -interaction=nonstopmode main.tex
bibtex main
pdflatex --disable-installer -interaction=nonstopmode main.tex
pdflatex --disable-installer -interaction=nonstopmode main.tex

docker compose up -d --build
docker compose exec backend python -m pytest
npm --prefix app/frontend/backoffice-app run test:unit
npm --prefix app/frontend/client-app run test:unit
npm --prefix app/frontend/client-app run test:e2e
\end{verbatim}
}
